\section{Praktikumsbericht}

\subsection{19.05.26 -- Dienstag}

Nach dem Gespräch habe ich mir bereits Gedanken darüber gemacht, wie das Frontend gestaltet werden soll.

\subsection{20.05.26 -- Mittwoch}

Ich habe zunächst alle vertraglichen Angelegenheiten bis 13 Uhr geregelt. Danach bekam ich die Anweisung, mir die Repositories im Raw-Format anzusehen, um die Header zu verstehen und ein besseres Gefühl dafür zu bekommen, wie alles funktioniert und zusammenhängt.
Dafür habe ich mir die wichtigsten Inhalte mit Unterstützung von KI erklären lassen, um schnell erste Zusammenhänge zu verstehen und offene Fragen zu klären.

\begin{itemize}
    \item STEMgraph abfragen
\end{itemize}

\subsection{21.05.26 -- Donnerstag}

Mein Auftrag war es, mich weiter zu vertiefen. Ich sollte mir bei \texttt{LinkedList} die Pointer \texttt{next\_ptr} und \texttt{prev\_ptr} in C ansehen und ein Verständnis dafür aufbauen. Dabei habe ich mit KI einige Punkte erarbeitet:

\begin{itemize}
    \item Was ich unter dem STEMgraph-Konzept verstehe.
    \item Was ich beim Lesen der Raws beachten sollte.
    \item Welchen Zusammenhang die Metadaten haben.
    \item Welche relevanten Punkte es bei einer LinkedList gibt.
    \item Welche Arten einer LinkedList es gibt.
    \item Wie eine Singly LinkedList und eine Doubly LinkedList aufgebaut sind.
    \item Welche Unterschiede sie haben.
    \item Welchen Vorteil sie gegenüber einer ArrayList haben.
\end{itemize}

Außerdem sollte ich mir ein Tutorial zu Neo4j ansehen und mir die Frage stellen, warum LinkedList und Neo4j sich ähneln.
Ich habe mir zwei bis drei Videos angeschaut und direkt festgestellt, dass der Name die Datei ist und der Array die Beziehung darstellt -- grob gesagt. Was Neo4j genau ist und wie es funktioniert, schaue ich mir am nächsten Tag weiter an.

\subsection{22.05.26 -- Freitag}

Ich sollte mit dem Bericht starten und beschreiben, womit ich mich die Woche beschäftigt hatte, welche relevanten Punkte mir aufgefallen sind und welche Fragen ich bisher noch nicht klären konnte.
Dazu sollte ich ein Repository erstellen und eine Datei im LaTeX-Format hineinschreiben, aber noch nichts ausführen, sondern mir zunächst nur die LaTeX-Syntax ansehen.
Beim weiteren Erfragen von Neo4j und LaTeX sind mir folgende Fragen aufgekommen:

\begin{itemize}
    \item Wie bekomme ich die Dateien in den Wert der Eigenschaften?
    \item Inwiefern ist es hilfreich, die Challenges in LaTeX zu schreiben, um sie in Neo4j aufzurufen?
\end{itemize}

\subsection{26.06.26 -- Dienstag}

Ich sollte meinen bisherigen Bericht in LaTeX mit Makefile schreiben. Wenn das erledigt ist, soll ich das Repository committen und einen Link schicken. Zusätzlich sollte ich, wenn ich fertig bin, ein Bash-Script schreiben, das ich einem Challenge-Repository gebe und das dann rekursiv alle weiteren Repositories herunterlädt und darauf aufbaut, was ich übergeben habe.
\begin{itemize}
    \item Um den Bericht in LaTeX zu bekommen, musste ich zunächst die LaTeX-Distribution MacTeX und die VS-Code-Extension LaTeX Workshop installieren und einrichten.
    \item Als das erledigt war, musste ich Homebrew, den Paketmanager für Linux und macOS, installieren.
    \item Um dann Makefile mit GNU Make zu bekommen, musste das Terminal berechtigt werden, damit ich \texttt{make} über Homebrew im Terminal installieren kann. Zusätzlich brauchte ich noch die VS-Code-Extension Makefile Tools.
    \item Teilweise hatte ich dann schon mit dem Umbau in LaTeX begonnen.
\end{itemize}
\subsection{27.06.26 -- Mittwoch}

Ich hatte damit begonnen, den gesamten Bericht in LaTeX zu formatieren und die Makefile zu befüllen.

\begin{itemize}
    \item Um nicht alle Befehle selbst eingeben zu müssen.
    \item Ich habe mich zusätzlich damit auseinandergesetzt, um mehr Verständnis dafür zu bekommen und es etwas schöner aussehen zu lassen.
\end{itemize}
